\documentclass[10pt,a4paper]{amsart}
\usepackage[margin=19mm]{geometry}
\usepackage{amsmath,amssymb,mathtools}
\usepackage[scr=rsfs,cal=euler]{mathalfa}
\usepackage{xfrac}
\usepackage[unicode,hidelinks,bookmarks=false]{hyperref}
\newcommand{\ie}{i.e.\ }
\newcommand{\cf}{cf.\ }
\newcommand{\resp}{respectively\ }
\newcommand{\Subsection}{\subsection}
\providecommand{\nobreakdash}{\nobreak\hbox{-}}
\setlength{\emergencystretch}{2em}
\multlinegap0pt
\usepackage{bm}
\usepackage{bbm}
\usepackage{dsfont}
\usepackage{xspace}
\usepackage{cancel}
\usepackage{mathtools}
\usepackage{amsthm}
\makeatletter
\@ifundefined{lemm}{\newtheorem{lemm}{Lemma}}{}
\makeatother
\newcommand{\NN}{\mathbb{N}} 
\newcommand{\RR}{\mathbb{R}}
\newcommand{\rfl}[2]{\left( #1 \, \vert \, #2 \right)}
\newcommand{\supp}{\mbox{\rm supp}}
\title[Measurable solutions of an alternative functional equation]{Measurable solutions of an alternative functional equation}
\author{P\'eter T\'oth}
\date{Source: arXiv:2508.17118v1 (CC BY 4.0)}
\begin{document}
\maketitle

\noindent\emph{Source and licence.} Measurable solutions of an alternative functional equation. Version arXiv:2508.17118v1 (CC BY 4.0), distributed under the Creative Commons Attribution 4.0 International licence, as recorded in the arXiv licence metadata for this exact version and shown on the abstract page https://arxiv.org/abs/2508.17118v1. Full rights notice: licenses/arxiv-2508.17118-CC-BY-4.0.txt.\par\smallskip

\noindent\textit{Editorial scope.} Counted unit: the Lemm whose source label is \texttt{lemm-locally\_constant} in the article (source0.tex lines 802--820, source label \texttt{lemm-locally\_constant}), delivered together with the article's own complete original proof of it (source0.tex lines 822--964, one \texttt{proof} environment of 143 lines). No mathematics is written, repaired or re-derived: statement and proof are the article's own text line for line. Required context retained uncounted: eq-maineq (lines 129--135). Every definition, display and label of those lines is retained unchanged. Omitted from this package and not redistributed: 1--128, 136--801, 821--821, 965--1798. Nothing inside a retained excerpt is omitted or altered: every retained body is delivered exactly as the article's lines are written, so no omission receipt of any kind accompanies this package. Citations: the retained excerpts carry no citation command at all, so no bibliography entry of the article is transcribed and no reference is redistributed. Reference adaptation: none was needed and none was made; every reference inside the delivered unit resolves inside it, so no unresolved reference or citation is produced. Printed slips kept literal: no printed word, symbol, hypothesis or number of the counted statement or of its proof was altered. Layout and class: the article's own class is \texttt{amsart} with packages including latexsym, amsmath, amsfonts, amsthm, amssymb, ifthen, enumerate, fontenc; the wrapper is the lane's pinned \texttt{amsart} editorial wrapper at 10pt on A4, which restates the theorem-like environments the retained text needs and adds the article's own macro definitions that the retained text needs (\texttt{\textbackslash NN}, \texttt{\textbackslash RR}, \texttt{\textbackslash rfl}, \texttt{\textbackslash supp}), with extra wrapper packages (bm, bbm, dsfont, xspace, cancel, mathtools). The delivered numbering is the unit's own. Assets: no retained span contains a figure environment, a graphics-inclusion command, a PDF-page inclusion, a table or a TikZ picture; no image file is copied into this package, the package declares no asset dependency and no image file is redistributed with it. Licence: version arXiv:2508.17118v1 of this article is distributed under the Creative Commons Attribution 4.0 International licence, as recorded in the arXiv licence metadata for this exact version and shown on the abstract page https://arxiv.org/abs/2508.17118v1. A full rights notice is delivered at licenses/arxiv-2508.17118-CC-BY-4.0.txt. Characters: every retained excerpt is pure ASCII, so the delivered engine, XeLaTeX with its default Latin Modern text font, reports no missing character.
\begin{equation}\label{eq-maineq} 
\varphi \left( \frac{x+y}{2} \right) 
\bigl( \psi_1(x) - \psi_2(y) \bigr) = 0 
\hspace{20mm} 
\left( \mbox{ for all } 
x \in I_1 \mbox{ and } y \in I_2 \right). 
\end{equation} 

\begin{lemm}\label{lemm-locally_constant}
Let 
$ \left( \varphi \,, \psi_1 \,, \psi_2 \right) $ 
be a solution of \eqref{eq-maineq}. 
Suppose that $ \varphi $ is measurable, 
and $ s \in J $ is such a point that 
$ \mu \left( \, \supp \varphi \cap 
\left[ s, s + \delta \right] \, \right) > 0 $ 
for every $ \delta > 0 $. 
Then, if $ \emptyset \neq K \subseteq I_1 $ 
is an interval fulfilling 
$ \rfl{K}{s} \subseteq I_2 \, $, then 
there exists $ \lambda \in \RR $ for which 
\[
\psi_1(x) = \psi_2(y) = \lambda 
\qquad \mbox{ holds for all } 
x \in K \mbox{ and } y \in \rfl{K}{s}. 
\]
\end{lemm}

\begin{proof}
Let 
$ x,y \in K $ be arbitrary numbers, 
fulfilling $ x < y $. 
Now 
\[ 
\rfl{ \, [x,y] \, }{s} = 
\left[ 2s-y, 2s-x \right] 
\subset \rfl{K}{s} \subseteq I_2 \,, 
\]
using the definition of the reflection operator, 
and the assumption for $ K $. 
Considering that $ I_2 $ is open, 
there exists $ d > 0 $ such that 
$ 2s - x + 2d \in I_2 \,$, so 
\[
\rfl{ \, [x,y] \, }{ \, [s, s+d] \, } = 
\left[ 2s-y , 2s + 2d -x \right] \subseteq I_2 
\]
holds. 
On the other hand, we have supposed 
$ \mu \left( \, \supp \varphi \cap 
\left[ s, s + d \right] \, \right) > 0 $. 
Therefore, applying the 
theorem of Steinhaus, the difference set 
\[
\Delta := \lbrace \, t_1 - t_2 \ : \ 
t_1 \,, t_2 \in \supp \varphi \cap [s,s+d] \, \rbrace
\] 
contains an open neighborhood of $ 0 $, 
i.e. there exists 
$ \varepsilon > 0 $ such that 
$ \left] - \varepsilon , \varepsilon 
\right[ \subseteq \Delta $. 
Now let $ m \in \NN $ be such that 
\[
h := \frac{y-x}{2m} \, \in \, 
\left] \, 0, \varepsilon \, \right[ 
\, \mbox{ is fulfilled}.
\]
Then there exist 
$ s_1 \,, s_2 \in \supp \varphi \cap [s, s+d] $ 
such that 
$ s_2 = s_1 + h $. 
Let us introduce new notations 
for some particular points: 
\[
x_j := x + j \cdot 2h \,, \qquad 
\widehat{x}_j := \rfl{x_j}{s_1} \ \mbox{ and } \ 
\widehat{\widehat{x}}_j := \rfl{\widehat{x}_j}{s_2} 
\qquad (j = 0,1, \dots , m).
\]
Obviously, 
$ x_0 = x $, $ x_m = y $ and $ x_j \in [x,y] $. 
Therefore 
$ \widehat{x}_j \in 
\rfl{\, [x,y] \,}{\, [s,s+d] \,} \subseteq I_2 $ 
for $ j=0,1, \dots ,m $ while 
\[
\widehat{\widehat{x}}_j = 2s_2 - \widehat{x}_j = 
2s_2 - 2s_1 + x_j = x_j + 2h = x_{j+1} \in I_1 
\qquad \mbox{for } j = 0,1, \dots , m-1. 
\] 
At this point let us utilize the fact that 
$ \left( \varphi \,, \psi_1 \,, \psi_2 \right) $ 
is a solution of \eqref{eq-maineq}. 
In particular, 
\begin{align*}
\varphi \left( 
\frac{x_j + \widehat{x}_j}{2} \right) 
\bigl( 
\psi_1 (x_j) - \psi_2(\widehat{x}_j) 
\bigr) &= 0 
\qquad (j = 0,1, \dots m) \ \mbox{ and } \\ 
\varphi \left( 
\frac{\widehat{\widehat{x}}_j + \widehat{x}_j}{2} \right) 
\bigl( 
\psi_1 (\widehat{\widehat{x}}_j) - \psi_2(\widehat{x}_j) 
\bigr) &= 0 
\qquad (j = 0,1, \dots m-1), 
\end{align*} 
which is equivalent to 
\begin{align*} 
\varphi \left( s_1 \right) 
\bigl( 
\psi_1 (x_j) - \psi_2(\widehat{x}_j) 
\bigr) &= 0 
\qquad (j = 0,1, \dots m) \ \mbox{ and } \\ 
\varphi \left( s_2 \right) 
\bigl( 
\psi_1 (x_{j+1}) - \psi_2(\widehat{x}_j) 
\bigr) &= 0 
\qquad (j = 0,1, \dots m-1).  
\end{align*} 
As $ s_1 \,, s_2 \in \supp \varphi $, 
these equations imply 
\[ 
\psi_1(x_j) = \psi_2( \widehat{x}_j ) = \psi_1(x_{j+1}) 
\qquad \mbox{ for every } 
j = 0,1, \dots ,m-1.
\]
Therefore 
$ \psi_1(x) = \psi_1 (x_0) = \psi_1 (x_1) = \dots = 
\psi_1 (x_m) = \psi_1(y) $ 
holds as well. 
Since $ x,y \in K $ were arbitrary, it follows 
that there exists 
$ \lambda \in \RR $ such that 
$ \psi_1(x) = \lambda$  for all $ x \in K $. 

Finally, we have to verify that 
$ \psi_2 $ is also constant on $ \rfl{K}{s} $ 
with value $ \lambda $. 
For this purpose let observe that 
$ \left( \varphi \,, \psi_1 \,, \psi_2 \right) $ 
is a solution of \eqref{eq-maineq} if, 
and only if, 
$ \left( \varphi \,, \psi_2 \,, \psi_1 \right) $ 
is a solution. 
Thus let us apply the same proof we have just carried out, 
but now for the solution 
$ \left( \varphi \,, \psi_2 \,, \psi_1 \right) $, 
and with the starting interval 
$ \rfl{K}{s} \subseteq I_2 \,$. 
Since 
$ \rfl{\rfl{K}{s}}{s} = K \subseteq I_1 \,$, 
the assumptions of the lemma are fulfilled 
for the solution 
$ \left( \varphi \,, \psi_2 \,, \psi_1 \right) $. 
Hence, according to the first part of our proof, 
there exists $ \nu \in \RR $ such that 
$ \psi_2(y) = \nu $ for all 
$ y \in \rfl{K}{s} $. 
But, for the 
points 
$ x_0 = x \in K $ and 
$ \widehat{x}_0 \in \rfl{K}{s}$ 
introduced before, we had 
\[
\lambda = \psi_1(x_0) = \psi_2(\widehat{x}_0) = \nu \,, 
\]
which completes the proof. 
\end{proof}

\end{document}
